% Part: normal-modal-logic
% Chapter: syntax-and-semantics
% Section: substitution
%READERNOTE{SOL6-A146: क्रमाने केलेल्या आदेशनाच्या दोन्ही मध्यवर्ती displays मध्ये संपूर्ण सूत्र दुसऱ्या आदेशनाच्या आत घेतले. पूर्वपक्षातील चलही बदलले पाहिजेत; केवळ उत्तरपक्ष बदलणे अपुरे आहे. अंतिम विस्तारित सूत्रे मूळातच योग्य होती.}

\documentclass[../../../include/open-logic-section]{subfiles}

\begin{document}

\olfileid{nml}{syn}{sub}

\olsection{एकाच वेळी केलेले आदेशन}

!!a{formula}~$!A$ मधील !!a{propositional variable} च्या सर्व
आस्थितींच्या जागी दुसरे !!{formula} घालून मिळणारे सूत्र हा
$!A$ चा \emph{दृष्टांत} असतो. वैधता आणि !!{derivability}
यांवर चर्चा करताना !!{formula}s च्या दृष्टांतांचा वारंवार
उल्लेख करावा लागेल. म्हणून ही संकल्पना नेमकी परिभाषित
करणे उपयुक्त आहे.

\begin{defn}\ollabel{def:subst-inst}
  $!A$ हे असे मोडल !!{formula} असू द्या की त्यातील सर्व
  !!{propositional
    variable}s ही $p_1$, \dots, $p_n$ यांपैकीच
  आहेत; तसेच $!D_1$, \dots, $!D_n$ हीदेखील मोडल !!{formula}s
  असू द्या. मग
  $\SSubst{!A}{\subst{!D_1}{p_1}, \dots, \subst{!D_n}{p_n}}$
  म्हणजे $!A$ मध्ये प्रत्येक $p_i$ च्या जागी $!D_i$ चे
  एकाच वेळी आदेशन करून मिळणारे सूत्र, अशी व्याख्या करतो.
  औपचारिकरीत्या ही $!A$ वरील विगमनाने केलेली व्याख्या आहे:
  \begin{enumerate}
    \tagitem{prvFalse}{\indcase{!A}{\lfalse}{%
        $\SSubst{\indfrm}{\subst{!D_1}{p_1}, \dots,
          \subst{!D_n}{p_n}}$ हे $\lfalse$ असते.}}{}
    \tagitem{prvTrue}{\indcase{!A}{\ltrue}{%
        $\SSubst{\indfrm}{\subst{!D_1}{p_1}, \dots,
          \subst{!D_n}{p_n}}$ हे $\ltrue$ असते.}}{}
    \item \indcase{!A}{q}{$\SSubst{\indfrm}{\subst{!D_1}{p_1}, \dots,
        \subst{!D_n}{p_n}}$ हे $q$ असते, जर $i
      = 1$, \dots,~$n$ साठी $q \not\ident p_i$ असेल.}
    \item \indcase{!A}{p_i}{$\SSubst{\indfrm}{\subst{!D_1}{p_1},
        \dots, \subst{!D_n}{p_n}}$ हे $!D_i$ असते.}
    \tagitem{prvNot}{\indcase{!A}{\lnot !B}{%
        $\SSubst{\indfrm}{\subst{!D_1}{p_1}, \dots, \subst{!D_n}{p_n}}$ हे
        $\lnot \SSubst{!B}{\subst{!D_1}{p_1}, \dots, \subst{!D_n}{p_n}}$ असते.}}{}
    \tagitem{prvAnd}{\indcase{!A}{(!B \land
        !C)}{$\SSubst{\indfrm}{\subst{!D_1}{p_1}, \dots,
          \subst{!D_n}{p_n}}$ हे पुढील सूत्र असते:
        \[(\SSubst{!B}{\subst{!D_1}{p_1}, \dots, \subst{!D_n}{p_n}} \land
          \SSubst{!C}{\subst{!D_1}{p_1}, \dots, \subst{!D_n}{p_n}}).\]}}{}
    \tagitem{prvOr}{\indcase{!A}{(!B \lor
          !C)}{$\SSubst{\indfrm}{\subst{!D_1}{p_1}, \dots,
          \subst{!D_n}{p_n}}$ हे पुढील सूत्र असते:
        \[(\SSubst{!B}{\subst{!D_1}{p_1}, \dots, \subst{!D_n}{p_n}} \lor
          \SSubst{!C}{\subst{!D_1}{p_1}, \dots, \subst{!D_n}{p_n}}).\]}}{}
    \tagitem{prvIf}{\indcase{!A}{(!B \lif
          !C)}{$\SSubst{\indfrm}{\subst{!D_1}{p_1}, \dots,
          \subst{!D_n}{p_n}}$ हे पुढील सूत्र असते:
        \[(\SSubst{!B}{\subst{!D_1}{p_1}, \dots, \subst{!D_n}{p_n}} \lif
          \SSubst{!C}{\subst{!D_1}{p_1}, \dots, \subst{!D_n}{p_n}}).\]}}{}
    \tagitem{prvIff}{\indcase{!A}{(!B \liff
          !C)}{$\SSubst{\indfrm}{\subst{!D_1}{p_1}, \dots,
          \subst{!D_n}{p_n}}$ हे पुढील सूत्र असते:
        \[(\SSubst{!B}{\subst{!D_1}{p_1}, \dots, \subst{!D_n}{p_n}} \liff
          \SSubst{!C}{\subst{!D_1}{p_1}, \dots, \subst{!D_n}{p_n}}).\]}}{}
    \item \indcase{!A}{\Box !B}{%
        $\SSubst{\indfrm}{\subst{!D_1}{p_1}, \dots, \subst{!D_n}{p_n}}$ म्हणजे
        $\Box
        \SSubst{!B}{\subst{!D_1}{p_1}, \dots, \subst{!D_n}{p_n}}.$}
    \tagitem{prvDiamond}{\indcase{!A}{\Diamond !B}{%
        $\SSubst{\indfrm}{\subst{!D_1}{p_1}, \dots, \subst{!D_n}{p_n}}$ म्हणजे
        $\Diamond
        \SSubst{!B}{\subst{!D_1}{p_1}, \dots, \subst{!D_n}{p_n}}.$}}{}
  \end{enumerate}
  !!{formula} $\SSubst{!A}{\subst{!D_1}{p_1}, \dots,
    \subst{!D_n}{p_n}}$ याला $!A$ चा \emph{आदेशन-दृष्टांत}
  म्हणतात.
\end{defn}

\begin{ex}
  समजा $!A$ हे $p_1 \lif \Box(p_1 \land p_2)$,
  $!D_1$ हे $\Diamond(p_2 \lif p_3)$ आणि $!D_2$ हे
  $\lnot\Box p_1$ आहे. मग
  $\SSubst{!A}{\subst{!D_1}{p_1}, \subst{!D_2}{p_2}}$ म्हणजे
  \begin{align*}
    \Diamond(p_2 \lif p_3) &
    \lif \Box(\Diamond(p_2 \lif p_3) \land \lnot\Box p_1)
    \intertext{तर $\SSubst{!A}{\subst{!D_2}{p_1}, \subst{!D_1}{p_2}}$ म्हणजे}
    \lnot\Box p_1 &
    \lif \Box(\lnot\Box p_1 \land \Diamond(p_2 \lif p_3))
    \intertext{लक्षात घ्या: एकाच वेळी केलेले आदेशन
      सर्वसाधारणपणे क्रमाने केलेल्या आदेशनासारखे नसते.
      उदाहरणार्थ, वरील
      $\SSubst{!A}{\subst{!D_1}{p_1}, \subst{!D_2}{p_2}}$ ची
      $\Subst{(\Subst{!A}{!D_1}{p_1})}{!D_2}{p_2}$ शी तुलना करा; नंतरचे असे आहे:}
    & \Subst{(\Diamond(p_2 \lif p_3) \lif
      \Box(\Diamond(p_2 \lif p_3) \land p_2))}{\lnot\Box p_1}{p_2},
      \text{ म्हणजे}\\
    \Diamond(\lnot\Box p_1 \lif p_3) &
    \lif \Box(\Diamond(\lnot\Box p_1
    \lif p_3) \land \lnot\Box p_1)\\
    \intertext{आणि $\Subst{(\Subst{!A}{!D_2}{p_2})}{!D_1}{p_1}$ शीही तुलना करा:}
    & \Subst{(p_1 \lif \Box(p_1 \land \lnot\Box p_1))}
      {\Diamond(p_2 \lif p_3)}{p_1}, \text{ म्हणजे}\\
    \Diamond(p_2 \lif p_3) &
    \lif \Box(\Diamond(p_2
    \lif p_3) \land \lnot\Box\Diamond(p_2 \lif p_3)).
  \end{align*}
\end{ex}

\end{document}
